%% Consumer-example — same content as report.md, rendered via LaTeX.
%% The Makefile sets TEXINPUTS=.:.templates: so \documentclass finds
%% kth-document.cls and \kthlogopath finds KTH_logo_RGB_bla.{pdf,png} in the
%% submodule without any path prefix in the .tex source.

\documentclass[english,11pt]{kth-document}

\title{Widget Library v2 --- User Guide}
\subtitle{Installing, using, and extending the widget library}
\doctype{Package documentation}
\author{Maintainers}
\affiliation{KTH Royal Institute of Technology}
\date{\today}
\version{2.0}
\shorttitle{Widget Library v2}

\begin{document}
\maketitle

\begin{kthbox}
\textbf{Summary.} This document describes how to install and use the
widget library in downstream projects. It is built from \texttt{report.tex}
using the \href{https://github.com/kth/kth-doc-templates}{kth-doc-templates}
LaTeX class, included here as a git submodule under \texttt{.templates/}.
\end{kthbox}

\section{Installation}

Add the package to your project:

\begin{verbatim}
npm install @kth/widget-library
\end{verbatim}

\section{Quick start}

Import the main entry point and instantiate a widget:

\begin{verbatim}
import { Widget } from '@kth/widget-library';

const w = new Widget({ kind: 'gauge' });
w.render('#root');
\end{verbatim}

\section{API reference}

\subsection{\texttt{new Widget(options)}}

Creates a new widget instance.

\begin{center}
\begin{tabular}{llll}
  \hline
  \textbf{Option} & \textbf{Type} & \textbf{Default} & \textbf{Description} \\
  \hline
  \texttt{kind}  & string & \texttt{'bar'} & Widget type --- bar, gauge, \ldots \\
  \texttt{theme} & string & \texttt{'kth'} & Visual theme. \\
  \texttt{data}  & array  & \texttt{[]}    & Initial dataset. \\
  \hline
\end{tabular}
\end{center}

\subsection{\texttt{widget.render(selector)}}

Mounts the widget into the DOM element matched by \texttt{selector}.

\begin{kthnotebox}
\textbf{Note.} The widget mounts synchronously and assumes the target
element already exists. Wait for \texttt{DOMContentLoaded} if you call
this from a top-level script.
\end{kthnotebox}

\section{Contact}

For support, contact the package maintainers at
\href{mailto:widget-library@kth.se}{widget-library@kth.se}.

\end{document}
